\documentclass[11pt]{article}
\usepackage[margin=1in]{geometry}
\usepackage[T1]{fontenc}
\usepackage[utf8]{inputenc}
\usepackage{amsmath,amssymb,amsthm}
\usepackage{enumitem}

\title{ICPC World Finals 2023\\B. Schedule}
\author{}
\date{}

\begin{document}
\maketitle

\section*{Problem Summary}

Each team contributes exactly one of its two members every week. For every pair of people from different
teams, we care about the largest gap between consecutive weeks when they meet, also counting the gaps
before the first meeting and after the last one. We must minimize the maximum such gap over $w$ weeks.

\section*{Key Observations}

\begin{itemize}[leftmargin=*]
    \item If a schedule has isolation $k$, then during every block of $k$ consecutive weeks each pair of
    people from different teams meets at least once.
    \item So it is enough to build a schedule of length $k$ where every such pair meets at least once. By
    repeating those $k$ weeks forever, we get the same isolation for any larger horizon.
    \item A team's schedule over $k$ weeks is a binary string of length $k$: bit $0$ means member $1$
    comes in, bit $1$ means member $2$ comes in.
    \item Two teams are compatible exactly when the four bit pairs $00,01,10,11$ all occur somewhere.
    Then every cross-team pair of employees meets at least once.
    \item If we take strings of length $k$ that start with $0$ and contain exactly $\lceil k/2\rceil$ ones,
    then any two distinct such strings are compatible. There are
    \[
    \binom{k-1}{\lceil k/2\rceil}
    \]
    such strings, and this is also the maximum possible number of pairwise compatible team strings.
\end{itemize}

\section*{Algorithm}

\begin{enumerate}[leftmargin=*]
    \item For $k=1,2,\dots,w$, find the first $k$ such that
    $\binom{k-1}{\lceil k/2\rceil} \ge n$.
    \item If no such $k$ exists, output \texttt{infinity}.
    \item Otherwise enumerate any $n$ distinct binary strings of length $k$ that start with $0$ and have
    exactly $\lceil k/2\rceil$ ones.
    \item Assign one string to each team.
    \item Output the first $w$ weeks by repeating these strings periodically with period $k$.
\end{enumerate}

\section*{Correctness Proof}

We prove that the algorithm returns the correct answer.

\paragraph{Lemma 1.}
If a schedule has isolation $k$, then the first $k$ weeks already contain at least one meeting for every
pair of people from different teams.

\paragraph{Proof.}
If some pair did not meet in the first $k$ weeks, then the initial gap before their first meeting would be at
least $k+1$, contradicting isolation $k$. \qed

\paragraph{Lemma 2.}
Any two distinct binary strings of length $k$ that start with $0$ and contain exactly $\lceil k/2\rceil$
ones are compatible.

\paragraph{Proof.}
Both strings start with $0$, so $00$ occurs in the first position. Since the strings are distinct but have the
same number of ones, there must be one position where they differ as $01$ and another where they differ
as $10$.

Among the last $k-1$ positions, each string has $\lceil k/2\rceil$ ones. Because
$2\lceil k/2\rceil > k-1$, the two sets of one-positions must intersect, so $11$ also occurs. Therefore all
four bit pairs occur, and the two teams are compatible. \qed

\paragraph{Lemma 3.}
No schedule block of length $k$ can contain more than $\binom{k-1}{\lceil k/2\rceil}$ pairwise compatible
team strings.

\paragraph{Proof.}
By flipping all bits of a team string if necessary, we may assume every string starts with $0$ without
changing compatibility.

Now remove the first bit. If one suffix were contained in another, then one of the pairs $01$ or $10$
would never occur, so the family of suffixes is an antichain in the Boolean lattice on $k-1$ elements.
Sperner's theorem gives the extremal size, and the standard complement-pair refinement in the odd case
reduces the bound exactly to $\binom{k-1}{\lceil k/2\rceil}$. \qed

\paragraph{Theorem.}
The algorithm outputs an optimal schedule whenever the answer is finite, and outputs
\texttt{infinity} otherwise.

\paragraph{Proof.}
Let $k$ be the first value found by the algorithm. By Lemma 2, the constructed $k$-week pattern makes
every pair of teams compatible, so every pair of employees from different teams meets at least once in
that block. Repeating the block makes each such pair meet at least once in every consecutive block of
$k$ weeks, so the isolation is at most $k$.

Conversely, if some schedule achieved smaller isolation, then by Lemma 1 its first block of that smaller
length would already contain $n$ pairwise compatible team strings, contradicting Lemma 3 and the
minimality of $k$. If no $k \le w$ satisfies the binomial bound, then no finite-isolation schedule exists.
\qed

\section*{Complexity Analysis}

Trying all $k \le w$ is $O(w)$. Generating the first $n$ valid patterns takes $O(nk)$ time, and printing
the repeated schedule takes $O(nw)$ time. The memory usage is $O(nk)$.

\section*{Implementation Notes}

\begin{itemize}[leftmargin=*]
    \item Binomial coefficients are computed with 64-bit arithmetic and capped at $n$, because larger
    values are irrelevant.
    \item Enumeration is done by backtracking with the first bit fixed to $0$.
\end{itemize}

\end{document}
